\chapter{Dlaczego nie ma ogólnego wzoru od stopnia piątego}

\section{Co znaczy „wzór ogólny”?}

Niech $a_1,\dots,a_n$ będą algebraicznie niezależnymi zmiennymi nad $\Q$. Ogólny wielomian unormowany stopnia $n$ to
\[
F(T)=T^n+a_1T^{n-1}+\cdots+a_n\in\Q(a_1,\dots,a_n)[T].
\]
Ogólny wzór przez pierwiastniki oznaczałby wyrażenie jego zer za pomocą $a_i$, działań $+,-,\cdot,/$ i skończenie wielu operacji wyciągania pierwiastków. Taki wzór można specjalizować do konkretnych współczynników, o ile nie zeruje się któryś z mianowników i dobierze się odpowiednie gałęzie. Algebraiczna niezależność oznacza, że współczynniki nie spełniają żadnego szczególnego związku wielomianowego. To jest sens słowa \emph{ogólny} (generyczny): po specjalizacji grupa Galois może się zmniejszyć, dlatego twierdzenie o przypadku generycznym nie unieważnia wzorów dla konkretnych wielomianów.

\section{Grupa Galois ogólnego wielomianu}

Weźmy algebraicznie niezależne zmienne $x_1,\dots,x_n$ i elementarne wielomiany symetryczne
\[
e_k=\sum_{1\le i_1<\cdots<i_k\le n}x_{i_1}\cdots x_{i_k},\qquad(1\le k\le n).
\]
Wówczas
\[
\prod_{i=1}^n(T-x_i)=T^n-e_1T^{n-1}+e_2T^{n-2}-\cdots+(-1)^ne_n.
\]

\begin{lemat}
Wielomiany $e_1,\dots,e_n$ są algebraicznie niezależne nad $\Q$. Każdy symetryczny wielomian w $x_1,\dots,x_n$ jest wielomianem w $e_k$, a każda symetryczna funkcja wymierna jest funkcją wymierną w $e_k$.
\end{lemat}
\begin{proof}
Porządkujemy jednomiany leksykograficznie, z $x_1>\cdots>x_n$. Wyraz wiodący $e_k$ to $x_1\cdots x_k$, więc wyraz wiodący $e_1^{b_1}\cdots e_n^{b_n}$ ma wektor wykładników
\[
(b_1+\cdots+b_n,\ b_2+\cdots+b_n,\ \dots,\ b_n).
\]
Różne wektory $(b_1,\dots,b_n)$ dają różne wektory wynikowe. W niezerowej kombinacji wielomianowej w $e_k$ największy z tych wyrazów wiodących nie może się więc skasować. Dowodzi to niezależności algebraicznej.

Niech symetryczny wielomian $P$ ma wyraz wiodący $c x_1^{a_1}\cdots x_n^{a_n}$. Symetria wymusza $a_1\ge\cdots\ge a_n$: gdyby $a_i<a_j$ dla $i<j$, zamiana $x_i,x_j$ dałaby większy leksykograficznie wyraz występujący w $P$. Odejmujemy od $P$ wielomian
\[
c\,e_1^{a_1-a_2}e_2^{a_2-a_3}\cdots e_{n-1}^{a_{n-1}-a_n}e_n^{a_n}.
\]
Ma on ten sam wyraz wiodący i ten sam stopień całkowity co ten wyraz. Reszta nadal jest symetryczna, ma mniejszy wyraz wiodący i stopień nie większy niż początkowe $P$. Istnieje tylko skończenie wiele jednomianów takiego stopnia, więc powtarzanie procesu kończy się zerem. To daje żądane przedstawienie wielomianowe.

Dla symetrycznej funkcji $r=p/q$ połóżmy $Q=\prod_{\sigma\in S_n}\sigma(q)$. To niezerowy symetryczny wielomian. Iloczyn $rQ$ też jest wielomianem, ponieważ $q$ jest jednym z czynników $Q$, i jest symetryczny. Zarówno $rQ$, jak i $Q$ są wielomianami w $e_k$, więc $r=(rQ)/Q\in\Q(e_1,\dots,e_n)$.
\end{proof}

\begin{wniosek}
Grupa Galois ogólnego wielomianu stopnia $n$ nad ciałem jego współczynników jest izomorficzna z $S_n$.
\end{wniosek}
\begin{proof}
Ciało $\Q(x_1,\dots,x_n)$ jest ciałem rozkładu wypisanego wielomianu nad $\Q(e_1,\dots,e_n)$. Każda permutacja niezależnych zmiennych $x_i$ daje automorfizm ciała funkcji wymiernych i ustala wszystkie $e_k$. Odwrotnie, automorfizm nad $\Q(e_1,\dots,e_n)$ musi permutować zera $x_i$, a te generują całe ciało. Grupą jest więc dokładnie $S_n$. Udowodniona niezależność algebraiczna $e_k$ pozwala utożsamić niezależne współczynniki $a_k$ z $(-1)^ke_k$. Przy tym utożsamieniu otrzymujemy właśnie ogólny wielomian $F$.
\end{proof}

\section{Nierozwiązalność $S_n$ dla $n\ge5$}

Nie potrzebujemy ogólnego twierdzenia o prostocie grup $A_n$. Wystarczy wykazać, że $A_n$ jest nietrywialna i równa własnej podgrupie pochodnej. Permutacje składamy od prawej do lewej.

\begin{lemat}
Dla $n\ge5$ zachodzi $A_n'=A_n\ne\{1\}$.
\end{lemat}
\begin{proof}
Najpierw $3$-cykle generują $A_n$. Każda permutacja jest iloczynem transpozycji, a parzystą można pogrupować w pary transpozycji. Para identycznych znika. Para z jednym wspólnym punktem jest $3$-cyklem lub jego odwrotnością. Dla czterech różnych punktów mamy
\[
(ab)(cd)=(acb)(acd).
\]
Zatem również para rozłącznych transpozycji jest iloczynem $3$-cykli.

Dla dowolnych trzech różnych punktów $a,b,c$ wybieramy jeszcze dwa różne punkty $d,e$. Permutacje $u=(ab)(de)$ i $v=(bc)(de)$ należą do $A_n$, są inwolucjami i $uv=(abc)$. Wobec tego
\[
[u,v]=uvuv=(abc)^2=(acb).
\]
Podgrupa pochodna zawiera odwrotności swoich elementów, więc zawiera także $(abc)$. Zawiera wszystkie $3$-cykle, a zatem całe $A_n$. W drugą stronę $A_n'\subseteq A_n$ zawsze. Grupa jest nietrywialna, bo zawiera $(123)$.
\end{proof}

\begin{wniosek}
Dla $n\ge5$ grupa $S_n$ nie jest rozwiązalna.
\end{wniosek}
\begin{proof}
Szereg pochodny $A_n$ pozostaje równy $A_n$, więc nie dochodzi do $1$. Gdyby $S_n$ była rozwiązalna, jej podgrupa $A_n$ też byłaby rozwiązalna, wbrew poprzedniemu lematowi.
\end{proof}

\section{Twierdzenie Abela--Ruffiniego}

\begin{twierdzenie}[Abela--Ruffiniego]
Dla każdego $n\ge5$ nie istnieje ogólny wzór przez pierwiastniki na zera wielomianu stopnia $n$ w zależności od jego współczynników.
\end{twierdzenie}
\begin{proof}
Grupą Galois ogólnego wielomianu jest $S_n$, na mocy wniosku 5.2. Dla $n\ge5$ grupa ta nie jest rozwiązalna (wniosek 5.4). Kryterium Galois (twierdzenie 4.8) mówi zatem, że ogólny wielomian nie jest rozwiązalny przez pierwiastniki.
\end{proof}

\section{Konkretny nierozwiązalny wielomian piątego stopnia}

Do następnego przykładu potrzebne są dwa krótkie fakty. Podajemy ich dowody, aby wniosek o grupie $S_5$ nie zależał od niewypowiedzianego testu nierozkładalności ani od klasyfikacji podgrup $S_5$.

\begin{lemat}[Kryterium Eisensteina dla wielomianu unormowanego]
Niech $f=T^n+a_{n-1}T^{n-1}+\cdots+a_0\in\Z[T]$. Jeśli pewna liczba pierwsza $p$ dzieli wszystkie $a_i$, lecz $p^2$ nie dzieli $a_0$, to $f$ jest nierozkładalny nad $\Q$.
\end{lemat}
\begin{proof}
Przypomnijmy potrzebną postać lematu Gaussa. Wielomian całkowity jest pierwotny, gdy największy wspólny dzielnik jego współczynników wynosi $1$. Iloczyn dwóch pierwotnych wielomianów jest pierwotny: gdyby pewne $p$ dzieliło wszystkie jego współczynniki, redukcje czynników byłyby niezerowymi wielomianami o zerowym iloczynie w $\F_p[T]$, co jest niemożliwe. Przez usunięcie mianowników i wyłączenie wspólnych dzielników każdy rozkład nad $\Q$ przyjmuje postać $f=cgh$, gdzie $g,h$ są pierwotne całkowite, a $c\in\Q^\times$. Jeśli także $f$ jest pierwotny całkowity, to $c=\pm1$: mianownik $c$ musiałby dzielić wszystkie współczynniki $gh$, a licznik byłby wspólnym dzielnikiem współczynników $f$. Dla unormowanego $f$ możemy więc przyjąć czynniki $g,h$ unormowane i całkowite.

Gdyby taki nietrywialny rozkład $f=gh$ istniał, redukcja modulo $p$ dawałaby $\overline g\,\overline h=T^n$. Oba czynniki są unormowane i zachowują stopnie, więc muszą być dodatnimi potęgami $T$. Ich wyrazy wolne byłyby podzielne przez $p$, a wyraz wolny $f$ przez $p^2$. To sprzeczność.
\end{proof}

\begin{lemat}[Cauchy, potrzebna postać]
Jeżeli liczba pierwsza $p$ dzieli rząd skończonej grupy $G$, to $G$ ma element rzędu $p$.
\end{lemat}
\begin{proof}
Rozważmy krotki $(g_1,\dots,g_p)$ spełniające $g_1\cdots g_p=1$. Jest ich $|G|^{p-1}$, bo ostatni element jest wyznaczony przez pozostałe. Przesunięcie cykliczne zachowuje ten zbiór: z $g_1\cdots g_p=1$ wynika $g_2\cdots g_pg_1=g_1^{-1}(g_1\cdots g_p)g_1=1$. Orbity przesunięcia mają rozmiar $1$ albo $p$. Stałe krotki mają postać $(g,\dots,g)$ z $g^p=1$. Ich liczba jest więc podzielna przez $p$, tak jak liczba wszystkich krotek. Jedna taka krotka odpowiada $g=1$; musi istnieć jeszcze inna. Dla jej elementu $g\ne1$ rząd wynosi $p$.
\end{proof}

\begin{twierdzenie}
Wielomian $f(x)=x^5-4x+2$ ma grupę Galois $S_5$ nad $\Q$, a więc nie jest rozwiązalny przez pierwiastniki.
\end{twierdzenie}
\begin{proof}
Kryterium Eisensteina (lemat 5.6) dla $p=2$ pokazuje nierozkładalność: wszystkie współczynniki poza wiodącym są podzielne przez $2$, a wyraz wolny nie jest podzielny przez $4$. Grupa Galois działa więc przechodnio na pięciu zerach. Z twierdzenia o orbicie i stabilizatorze jej rząd jest podzielny przez $5$. Lemat 5.7 daje element rzędu $5$; w $S_5$ musi być on $5$-cyklem, ponieważ rząd permutacji jest najmniejszą wspólną wielokrotnością długości jej rozłącznych cykli.

Pochodna $f'(x)=5x^4-4$ ma dwa rzeczywiste punkty krytyczne $\pm a$, gdzie $a=(4/5)^{1/4}$. Funkcja rośnie na $(-\infty,-a)$, maleje na $(-a,a)$ i rośnie na $(a,\infty)$. Ponadto
\[
f(-a)=2+\frac{16}{5}a>0,\qquad
f(a)=2-\frac{16}{5}a<0.
\]
Ostatnia nierówność wynika z $a>5/8$, bo $a^4=4/5>(5/8)^4$. Granice $f$ w $-\infty$ i $+\infty$ wynoszą odpowiednio $-\infty$ i $+\infty$. Twierdzenie o wartości pośredniej oraz ścisła monotoniczność dają dokładnie po jednym zerze na każdym z trzech przedziałów. Zatem $f$ ma dokładnie trzy zera rzeczywiste i jedną parę nierzeczywistych zer sprzężonych. Sprzężenie zespolone zachowuje ciało rozkładu i działa na zerach jako pojedyncza transpozycja.

Oznaczmy punkty cyklu przez $\Z/5\Z$ tak, by $5$-cykl był przesunięciem $k\mapsto k+1$. Jeżeli transpozycja zamienia $u$ i $v$, to jej sprzężenia potęgami cyklu zamieniają $u+k$ i $v+k$. Krawędzie te tworzą spójny graf, ponieważ $u-v\ne0$ generuje addytywną grupę $\Z/5\Z$. Transpozycje krawędzi spójnego grafu generują wszystkie transpozycje: wzdłuż prostej ścieżki można kolejno otrzymywać transpozycję jej końców, używając tożsamości $(ab)(bc)(ab)=(ac)$. Wszystkie transpozycje generują grupę symetryczną. Zatem grupa Galois jest $S_5$.
\end{proof}

\section{Czego twierdzenie nie mówi}
\begin{itemize}
\item Wielomiany $x^n-a$ są rozwiązalne przez pierwiastniki dla każdego $n$. Stopień sam w sobie nie przesądza o konkretnym równaniu.
\item Równanie $x^5-1=0$ ma zera $\zeta_5^j$. Jego grupa Galois nad $\Q$ jest abelowa, izomorficzna z $(\Z/5\Z)^\times$.
\item Nierozwiązalność przez pierwiastniki nie oznacza braku metod numerycznych ani braku wzorów używających funkcji specjalnych.
\item Twierdzenie dotyczy skończonych wyrażeń algebraicznych. Nie zabrania szeregów, granic, całek ani rozwiązań przybliżonych.
\end{itemize}


